\chapter{Lý sự cùn sát lớp thật}
\markboth{Lý sự cùn sát lớp thật}{Lý sự cùn sát lớp thật}

\section{Em không cãi, em chỉ không phục}
\begin{truyen}{Em không cãi, em chỉ không phục}{Classroom Talk}
\chuhoa{H}{ọc sinh: ``Em không cãi đâu. Em chỉ không phục cách thầy nói thôi.''}
\textit{Student: ``I am not arguing. I just do not agree with the way you said it.''}\\[4pt]

\teacherpair{Không phục thì em có quyền nói. Nhưng nói để hiểu ra khác với nói để giành phần thắng.}{If you disagree, you have the right to say so. But speaking to understand is different from speaking to win.}

\studentpair{Thì em đang nói đó.}{Well, I am saying it right now.}

\teacherpair{Ừ. Và cả lớp đang nghe xem em muốn làm rõ vấn đề hay chỉ muốn giữ sĩ diện.}{Yes. And the whole class is listening to see whether you want to clarify the issue or simply protect your pride.}
\end{truyen}

\begin{baihoc}
Không phải cứ to tiếng là có lý. Nhiều khi chỉ là cái tôi đang bị dồn góc.
\textit{(A louder tone does not prove stronger logic. Sometimes it only reveals a cornered ego.)}
\end{baihoc}

\begin{ghinhoanh}
\textbf{Short English to remember:}

Disagreement is allowed.

Ego is often louder than reason.
\end{ghinhoanh}

\begin{tuhocnhanh}
1. Khi em phản ứng mạnh, em đang bảo vệ lẽ phải hay sĩ diện? \textit{(When you react strongly, are you defending truth or pride?)}

2. Em có biết nói khác ý mà vẫn giữ tôn trọng không? \textit{(Can you disagree while staying respectful?)}
\end{tuhocnhanh}
\ngancach

\section{Em ngủ nhưng đầu óc em vẫn online}
\begin{truyen}{Em ngủ nhưng đầu óc em vẫn online}{Classroom Talk}
\chuhoa{H}{ọc sinh: ``Mắt em nghỉ thôi, chứ đầu óc em vẫn online.''}
\textit{Student: ``Only my eyes were resting. My brain was still online.''}\\[4pt]

\teacherpair{Online vậy thầy vừa giảng tới đâu?}{If your brain was online, then where exactly did my lesson just reach?}

\studentpair{Dạ... tới phần em cần tỉnh rồi.}{Uh... it reached the part where I need to wake up.}

\teacherpair{Đúng. Phần đó bắt đầu từ ba phút trước.}{Exactly. That part started three minutes ago.}
\end{truyen}

\begin{baihoc}
Nói vui có thể cứu một pha quê. Nhưng nó không cứu được lỗ hổng vừa xảy ra.
\textit{(A joke may save a moment of embarrassment, but it does not repair the loss that just happened.)}
\end{baihoc}

\begin{ghinhoanh}
\textbf{Short English to remember:}

Humor can lighten the moment.

It cannot replace attention.
\end{ghinhoanh}

\begin{tuhocnhanh}
1. Em đang thiếu ngủ thật hay thiếu kỷ luật ngủ nghỉ? \textit{(Are you truly sleep-deprived, or lacking sleep discipline?)}

2. Em cần sửa gì để vào lớp tỉnh hơn? \textit{(What needs to change so you enter class more awake?)}
\end{tuhocnhanh}
\ngancach

\section{Em không quay bài, em chỉ nhìn cho đỡ hoảng}
\begin{truyen}{Em không quay bài, em chỉ nhìn cho đỡ hoảng}{Classroom Talk}
\chuhoa{H}{ọc sinh: ``Em không quay bài. Em chỉ liếc qua cho đỡ panic thôi.''}
\textit{Student: ``I was not cheating. I only glanced over so I would not panic.''}\\[4pt]

\teacherpair{Vậy nếu em chỉ lấy tiền người ta cho đỡ hoảng, có gọi là lấy không?}{Then if you took someone else’s money just to calm yourself down, would it still count as taking it?}

\studentpair{Thầy ví mạnh tay quá.}{That comparison is pretty harsh, sir.}

\teacherpair{Vì em đang làm nhẹ một chuyện không nhẹ.}{Because you are trying to make something serious sound harmless.}
\end{truyen}

\begin{baihoc}
Đổi tên một hành vi không làm hành vi đó sạch hơn.
\textit{(Renaming an action does not make it cleaner.)}
\end{baihoc}

\begin{ghinhoanh}
\textbf{Short English to remember:}

A softer label does not change the act.

Cheating remains cheating.
\end{ghinhoanh}

\begin{tuhocnhanh}
1. Em có hay làm nhẹ lỗi bằng cách đổi cách gọi không? \textit{(Do you soften your mistakes by renaming them?)}

2. Khi hoảng, em có cách nào khác ngoài nhìn bài người khác? \textit{(When you panic, what can you do besides looking at someone else’s paper?)}
\end{tuhocnhanh}
\ngancach

\section{Em không hỗn, em nói thẳng thôi}
\begin{truyen}{Em không hỗn, em nói thẳng thôi}{Classroom Talk}
\chuhoa{H}{ọc sinh: ``Em không hỗn. Em chỉ nói thẳng. Người lớn cứ thích vòng vo nên nghe nói thẳng mới khó chịu.''}
\textit{Student: ``I am not being disrespectful. I am just being direct. Adults like to speak in circles, so they get upset when they hear direct words.''}\\[4pt]

\teacherpair{Nói thẳng là nói đúng và rõ. Còn đâm người khác bằng lời rồi gọi đó là thẳng tính thì dễ quá.}{Being direct means speaking clearly and truthfully. Stabbing someone with words and then calling it honesty is far too easy.}

\studentpair{Vậy sao phân biệt?}{So how do you tell the difference?}

\teacherpair{Nói xong mà vấn đề sáng hơn, đó là thẳng. Nói xong chỉ làm người khác nghẹn lại, thường là thiếu tôn trọng.}{If the issue becomes clearer after you speak, that is directness. If all you do is leave someone feeling hurt and shut down, that is usually disrespect.}
\end{truyen}

\begin{baihoc}
Thẳng tính không phải giấy phép để thiếu lịch sự.
\textit{(Being direct is not a license for disrespect.)}
\end{baihoc}

\begin{ghinhoanh}
\textbf{Short English to remember:}

Direct is not the same as rude.

Clarity should not crush respect.
\end{ghinhoanh}

\begin{tuhocnhanh}
1. Em có đang dùng chữ ``thẳng'' để che một cách nói làm đau người khác không? \textit{(Do you use the word direct to hide hurtful speech?)}

2. Em có thể nói cùng ý đó theo cách nào tốt hơn? \textit{(How could you say the same thing in a better way?)}
\end{tuhocnhanh}
\ngancach